\section{The arithmetic square, Riemann--Roch, and absolute counting}
\label{sec:sites-twistor}

\subsection{Source inventory and the exact numerical conclusion}

The statements in this subsection distinguish a proved identity on test
functions from an algebro-geometric realization of that identity.  The
attached version of Connes--Consani, \emph{The Riemann--Roch strategy}
(arXiv:1805.10501), \S3.1, equations labelled \texttt{negcrit},
\texttt{negcrit1}, and \texttt{negcrit1bis}, gives the numerical target.
Its \S3.4 lists five steps: construction of a complex lift, regularization
of the relative intersection theory, a surface Riemann--Roch formula,
an existence theorem for sections, and tropical descent.  Construction of
the complex lift and a lift of the scaling correspondences is part of that
paper; the complete surface intersection/Riemann--Roch theorem is not.

Here are the relevant already proved results, with locators in the supplied
sources.  They are background, not findings of this work.
\begin{itemize}
\item \emph{Geometry of the scaling site} (1603.03191), Theorems 3.9
(\texttt{scaltop}), 4.2 (\texttt{structure2}), and 4.3
(\texttt{structure3}): points, stalks, and points over the tropical
semifield; Theorem 5.6 (\texttt{thmjaccp}): the divisor-class description
on the periodic orbit $C_p$; Theorem 5.17 (\texttt{RRperiodic}):
$\operatorname{Dim}_{\mathbb R}H^0(D)-\operatorname{Dim}_{\mathbb R}H^0(-D)
=\deg D$, with the continuous dimension defined as the limit in the
source's \texttt{rr1}.  These statements do not assume RH.
\item \emph{Riemann--Roch for the Arakelov compactification of
$\operatorname{Spec}\mathbb Z$} (2205.01391), Theorem 4.3
(\texttt{rrspzb}):
\[
 \dim_{\mathbb S[\pm1]}H^0(D)-\dim_{\mathbb S[\pm1]}H^1(D)
 =\left\lceil\frac{\deg D+\log2}{\log3}\right\rceil'
 -\mathbf1_L(\deg D).
\]
Here the odd ceiling and the exceptional set $L$ are exactly those in
\texttt{ceilingprime} and \texttt{excL} of that paper.  Theorem 5.3
(\texttt{serredual}) identifies
$H^0(K-D)$ with the appropriate tolerant-module internal Hom from
$H^1(D)$, where $K=-2\{2\}$.  Both statements are unconditional.
\item \emph{Riemann--Roch for the ring $\mathbb Z$} (2306.00456),
Theorem 5.1 (\texttt{rrspzb}):
\[
 \dim_{\mathbb S}H^0(D)-\dim_{\mathbb S}H^1(D)
 =c(\deg D/\log2)+1.
\]
The function $c$ is the source's \emph{right-continuous} extension of
$\lceil x\rceil$ for positive nonintegral $x$ and
$-\lceil-x\rceil$ for negative nonintegral $x$.  It must not be replaced
by the ordinary ceiling at integral arguments; for example $c(0)=1$.
The exact formula, including this convention, is already unconditional.
\end{itemize}
None is a Riemann--Roch theorem for the two-dimensional arithmetic square.
Their existence therefore supplies no missing surface theorem merely by
substitution of A.

Write $d^*u=du/u$ and
$Mf(s)=\int_0^\infty f(u)u^s\,d^*u$.  All test functions below belong to
$\mathcal D=C_c^\infty(\mathbb R_+^*,\mathbb C)$.
Put $\widetilde f(u)=u^{-1}f(u^{-1})$, and let $N$ be the distribution
of \texttt{negcrit1} in the preceding source:
\[
 N(h)=\sum_{n\geq1}\Lambda(n)h(n)
 +\int_1^\infty\frac{u^2h(u)-h(1)}{u^2-1}\,d^*u
 +\tfrac12(\log\pi+\gamma)h(1).
\]
The subtraction makes the integrand bounded at $u=1$ and the tail
integrable.  Thus this expression is defined on $\mathcal D$.
For complex functions the Hermitian extension of the real intersection
form is
\[
 I(f,g)=\tfrac12\left\{N(f*\widetilde{\overline g})+
             \overline{N(g*\widetilde{\overline f})}\right\}.
\]
This symmetrization matters: the one-sided distribution $N$ is not itself
an invariant trace on the whole multiplicative group.

\begin{proposition}[AC-G01: the full numerical Hodge inequality]
\label{AC-G01}
Under A, with the nontrivial zeta zeros written
$\rho=\tfrac12+i\gamma$ and $m_\rho$ their multiplicities,
\begin{equation}\label{eq:G01}
 2I(f,g)=Mf(0)\overline{Mg(1)}+Mf(1)\overline{Mg(0)}
 -\sum_\rho m_\rho Mf(\rho)\overline{Mg(\rho)}.
\end{equation}
The sum converges absolutely.  In particular, on
$\mathcal D_{00}=\{f:Mf(0)=Mf(1)=0\}$ the matrix
$\bigl(I(f_i,f_j)\bigr)_{i,j=1}^r$ is negative semidefinite for every finite
family.  For real $f$, this is precisely the inequality
$N(f*\widetilde f)\leq0$ requested in the Riemann--Roch strategy.
Equality on $\mathcal D_{00}$ holds exactly when $Mf(\rho)=0$ for every
nontrivial zeta zero.
\end{proposition}
\begin{proof}
The usual explicit formula in the normalization of $N$ is
\[
 N(h)+N(\widetilde h)=Mh(0)+Mh(1)-\sum_\rho m_\rho Mh(\rho).
\]
This is the same explicit formula used in the source's \S3.1, and in
CCM (math/0703392), \S7, Proposition \texttt{PosTrRHint} and equation
\texttt{terms01Tr}.  To check its normalization, the Mellin transform
of the distribution on $[1,\infty)$ is
$-\Lambda_\zeta'(s)/\Lambda_\zeta(s)$ for
$\Lambda_\zeta(s)=\pi^{-s/2}\Gamma(s/2)\zeta(s)$; adding its reflected
contour picks up the poles at $0,1$ and the negatives of the zero residues.
Thus there is no extra factor of two in this displayed identity.
Now
$M(f*\widetilde{\overline g})(s)=Mf(s)\overline{Mg(1-\overline s)}$
and
$\widetilde{(f*\widetilde{\overline g})}
 =\overline{g*\widetilde{\overline f}}$.
A makes $1-\overline\rho=\rho$, which gives \eqref{eq:G01}.
For every $B>0$, integration by parts after $u=e^x$ gives
$Mf(\tfrac12+it)=O_{f,B}((1+|t|)^{-B})$.
The zero-counting bound $N_\zeta(T)=O(T\log(T+2))$ recalled in
Connes 1999, \S I, equations (1)--(3) and the following paragraph,
proves absolute convergence.  On $\mathcal D_{00}$,
$-2I(f,f)=\sum m_\rho|Mf(\rho)|^2$.  Apply this identity to
$\sum c_if_i$ for the matrix assertion, and use that each summand is
nonnegative for the equality criterion.  If $f$ is real,
$f*\widetilde f$ is fixed by the tilde operation, so
$I(f,f)=N(f*\widetilde f)$.
\end{proof}

\begin{proposition}[AC-G02: a constructed numerical square]
\label{AC-G02}
Let $\mathcal H$ be the closure of the vectors
$\Phi(f)=(\sqrt{m_\rho}Mf(\rho))_\rho$, $f\in\mathcal D$,
in $\ell^2$ of the distinct zeros.  Under A the Hilbert space
$\mathcal H$ is well-defined.  The space
$\mathcal N=\mathbb C^2\oplus\mathcal H$, with Hermitian form
\[
 B((a,b,v),(a',b',v'))=a\overline{b'}+b\overline{a'}-
 \langle v,v'\rangle,
\]
contains the explicitly defined classes
$D_f=(Mf(0),Mf(1),\Phi(f))$, and
$B(D_f,D_g)=2I(f,g)$.
Its two coordinate classes $E=(1,0,0)$, $F=(0,1,0)$ satisfy
$E^2=F^2=0$, $B(E,F)=1$; on the orthogonal of $E+F$ the form
is negative definite.  Its minimal spectral Hilbert realization is unique
up to the isometry carrying each marked vector $\Phi(f)$ to its counterpart.
\end{proposition}
\begin{proof}
Square summability is AC-G01.  The form and the classes are given by the
displayed formulas, so their compatibility is \eqref{eq:G01}.
Orthogonality to $E+F$ means $a+b=0$; consequently
$B((a,-a,v),(a,-a,v))=-2|a|^2-\|v\|^2$, which is negative
unless the vector is zero.  For uniqueness, two dense families with the
same Gram kernel have an isometry on their finite linear spans:
$\sum c_i\Phi(f_i)\mapsto\sum c_i\Phi'(f_i)$.
Equality of squared norms proves that this map is well-defined, and its
continuous extension is onto by density.
\end{proof}

This constructs a numerical intersection space, not a surface, a category
of divisors, an integral N\'eron--Severi lattice, or a sheaf cohomology.
It forces any proposed trace-compatible geometric realization to have the
same negative form on the marked primitive correspondences.  In particular,
a model that gives one of these correspondences positive square is excluded.
It does not determine a geometric square: even at the numerical level one
may adjoin any extra negative definite space with zero degree and codegree,
leaving all the displayed marked intersections unchanged.  The quotient
by the marked subspace can have arbitrarily prescribed finite dimension.
Thus the missing construction remains the surface intersection/Riemann--Roch
and effectivity input in \S3.4, steps 2--5, of 1805.10501.  This is a
statement about the precise missing input, not a formal independence theorem.

\begin{proposition}[AC-G03: a finite single-Frobenius obstruction]
\label{AC-G03}
No finite collection of finite-dimensional vector spaces with invertible
operators $F_i$ and one fixed real $q>1$ can express the nontrivial-zero
divisor of $\zeta(s)$ as the divisor of
$\prod_i\det(1-q^{-s}F_i)^{n_i}$, $n_i\in\mathbb Z$.
\end{proposition}
\begin{proof}
Each nonzero eigenvalue $\alpha$ contributes only the arithmetic progression
$s=(\log\alpha+2\pi i k)/\log q$, $k\in\mathbb Z$,
with multiplicity bounded by the fixed dimension.  A finite product has
$O(T)$ zeros and poles, counted with multiplicity, in $|\Im s|\leq T$;
cancellation only decreases this bound.  The Riemann--von Mangoldt formula (Connes 1999, \S I, (2)--(3)) gives
$N_\zeta(T)=(T/2\pi)\log(T/2\pi)-T/2\pi+O(\log(T+2))$
for the upper half-plane, a contradiction.  This does not exclude
$\operatorname{Spec}\mathbb Z$ as a scheme or its usual prime-indexed Euler
product.  It excludes the stated finite, single-Frobenius determinant model.
\end{proof}

\subsection{The absolute twistor line and absolute cyclotomy}

In 2609.00299, Proposition 3.2 (\texttt{prop:points\_X\_K}) describes the
unfolded line over a field containing square roots of $-1$; Theorem 5.1
(\texttt{thm:twistor\_space\_scheme}) gives
$\sigma(z)=-1/\overline z$ on the folded complex line.  The absence of real
fixed points follows already from $|z|^2=-1$, with $0,\infty$ interchanged.
Proposition 7.3 (\texttt{prop:global\_sections\_frobenius}) states the
$\mathbb F_{1^2}$ global sections in its specified odd arithmetic topos.
Proposition 7.5 (\texttt{prop:twistor\_frobenius}) gives, for odd $n$,
\[
 f_n(z)=\begin{cases}z^n&n\equiv1\pmod4,\\-z^{-n}&n\equiv3\pmod4.
 \end{cases}
\]
These are source constructions independent of A.  Their direct coordinate
formulas obey $f_mf_n=f_{mn}$ and commute with $\sigma$: for
$h(z)=-z^{-1}$, all odd power maps commute with $h$, $h^2=1$, and
$f_n=h^{(1-\chi_{-4}(n))/2}\circ(z\mapsto z^n)$.
No zero-location assertion is involved in these verifications.

Likewise, 2112.08820, Theorem 2.11 (\texttt{w0s1}), identifies the
endomorphism invariant ring $W_0(\overline{\mathbb S})$ with
$\mathbb Z[\mathbb Q/\mathbb Z]$; Theorem 2.12
(\texttt{rabi10}) identifies the Frobenius/Verschiebung structure with
the integral BC-system.  These are unconditional statements in
endomorphism $K$-theory.  A supplies the zero-location theorem for the
Dirichlet factors that occur in their arithmetic realizations; it neither
changes these rings nor adds an intersection theorem to them.

The phrase ``Gaussian avatar'' has a precise safe arithmetic realization:
the Artin motive of $\operatorname{Spec}\mathbb Q(i)$, or the ordinary
scheme $\operatorname{Spec}\mathbb Z[i]$.  Its zeta is
$\zeta(s)L(s,\chi_{-4})$.  The presence of the same quadratic sign in the
absolute twistor formulas is not by itself a theorem that the zeta of the
entire absolute twistor object equals this Dedekind zeta.  Such an
identification requires a specified counting/trace formalism for that object.

\subsection{Artin--Tate objects, global zeros, and the prime 2}

Let $\varepsilon$ denote the nontrivial one-dimensional rational
representation of $C_2=\operatorname{Gal}(\mathbb Q(i)/\mathbb Q)$, and
let $\mathbb L$ denote the Tate object with good-prime Frobenius $p$.
Finite direct sums of $\mathbb L^j$ and $\varepsilon\mathbb L^j$,
$j\in\mathbb Z$, form the explicitly constructed tensor category of
finite-dimensional, integer-graded rational $C_2$-representations.
Its Grothendieck ring is
\[
 R=\mathbb Z[\mathbb L,\mathbb L^{-1},\varepsilon]/(\varepsilon^2-1).
\]
These are actual Artin--Tate motives over $\mathbb Q$: the sign summand
is cut out of $h^0(\operatorname{Spec}\mathbb Q(i))$ by the projector
$(1-c)/2$, and the other summand by $(1+c)/2$.
Over $U=\operatorname{Spec}\mathbb Z[1/2]$ their Artin part comes from
the finite \emph{\'etale} double cover $\operatorname{Spec}\mathbb Z[i,1/2]$.
The associated $\ell$-adic sheaves are the trivial and sign sheaves,
tensored with $\mathbb Q_\ell(-j)$ on the open set where $\ell$ is
invertible.  This specifies the objects and their realizations, rather
than only a polynomial that might be their trace.

\begin{proposition}[AC-G04: the complete sign--Tate construction]
\label{AC-G04}
For integers $a_j^+,a_j^-$ of finite support, the Artin--Tate class
$V=\sum_j(a_j^++a_j^-\varepsilon)\mathbb L^j$ has complete finite-prime
$L$-function
\begin{equation}\label{eq:G04}
 L(V,s)=\prod_j\zeta(s-j)^{a_j^+}L(s-j,\chi_{-4})^{a_j^-}.
\end{equation}
At an odd prime its $n$th Frobenius trace is
$\sum_j(a_j^++a_j^-\chi_{-4}(p)^n)p^{nj}$.
Under A the divisor contributed by the nontrivial zeros of the primitive
factors in \eqref{eq:G04} is supported on
$\Re s=j+\tfrac12$.  Negative multiplicities can produce poles and
cancellations; this is not a claim that a virtual class has only zeros.
Local purity is unconditional: every constituent in degree $j$ has
weight $2j$, since its good-prime eigenvalue is $p^j$ or $-p^j$.
A direct sum with distinct $j$ is graded by these weights, not pure of
one weight.
\end{proposition}
\begin{proof}
For $\Re s$ sufficiently large, use
$\exp(\sum_{n\geq1}\alpha^n p^{-ns}/n)=(1-\alpha p^{-s})^{-1}$.
At odd $p$ this gives
$\prod_j(1-p^{j-s})^{-a_j^+}
(1-\chi_{-4}(p)p^{j-s})^{-a_j^-}$.
At $2$, the sign representation has no inertia invariants and the
trivial one has one-dimensional invariants, so the factor is
$\prod_j(1-2^{j-s})^{-a_j^+}$.
Multiplication gives \eqref{eq:G04} in an absolutely convergent half-plane
and hence by meromorphic continuation.  The trace formula follows from
the stated eigenvalues.  Under A a nontrivial factor zero is
$s=j+\rho$ with $\Re\rho=1/2$; the divisor of a product is the sum of
its divisors.  Finally $|\pm p^j|=p^{(2j)/2}$ proves local purity without
using A.  Neither this computation nor A asserts one common weight for
several different $j$.
\end{proof}

\begin{proposition}[AC-G05: tensor multiplicativity and its ramified failure]
\label{AC-G05}
The assignment of local rational zeta functions of $V$ is a ring map
from $R$ into the product of rational Witt rings over odd primes,
where direct sum is Witt addition and tensor product is Witt multiplication.
It is injective.  Including the factor at $2$ defined by inertia
invariants destroys tensor multiplicativity.  Specifically,
\[
 L_2(\varepsilon,s)=1,
 \qquad L_2(\varepsilon\otimes\varepsilon,s)
 =(1-2^{-s})^{-1}.
\]
Thus the complete Euler-product assignment cannot be a ring map for this
primewise Witt multiplication.
\end{proposition}
\begin{proof}
At an odd prime the eigenvalues of a direct sum are their multiset union
and those of a tensor product are all pairwise products, which are the
stated Witt operations.  At $p=3$, the distinct eigenvalues $3^j$ and
$-3^j$ give distinct linear factors in $\mathbb Q(T)$.
Unique factorization therefore recovers every $a_j^+$ and $a_j^-$ from
that one local rational function, proving injectivity even for virtual
classes.  At $2$, inertia acts on $\varepsilon$ by $-1$, giving zero
invariants, whereas its square is the trivial representation with one
invariant.  Witt multiplication of the two unit rational functions
$1$ is still $1$, which contradicts the displayed tensor factor.
This is precisely the failure of
$(V\otimes W)^I=V^I\otimes W^I$ in this example.
\end{proof}

Over $U$ the removed factor also matters analytically:
$\zeta^{(2)}(s-j)=(1-2^{j-s})\zeta(s-j)$ has additional elementary
zeros on $\Re s=j$.  These are not nontrivial zeros of its primitive
completion and must not be relabelled violations of A or assigned the
line $j+1/2$.

\begin{proposition}[AC-G06: counts, categorical positivity, and the Tate tower]
\label{AC-G06}
In the category just constructed, a class is represented by an actual
object exactly when all $a_j^\pm\geq0$.  Nonnegative Frobenius trace
counts do not suffice: $1-\varepsilon$ has all good-prime power traces
$1-\chi_{-4}(p)^n\in\{0,2\}$, but is not the class of an object.
For every integer $k\geq0$, the object
$40h((\mathbb P^1)^{2k})(k)$ has zeta
\[
 Z_k(s)=\prod_{a=-k}^{k}\zeta(s-a)^{40\binom{2k}{k+a}}.
\]
A gives its nontrivial critical lines $\Re s=a+1/2$.
This constructs one object per word-count of the stated form; it does
not construct the multiplication or associativity constraints of a
categorification of all GCT structure constants.
\end{proposition}
\begin{proof}
The simple objects of the graded representation category are exactly
$\mathbb L^j,\varepsilon\mathbb L^j$, so their multiplicities in an actual
object are nonnegative and conversely determine such a direct sum.
The virtual class $1-\varepsilon$ has the displayed traces, but its
ordinary dimension is zero and it is nonzero in the free additive basis
of $R$, so it cannot be an actual object.
The projective-line decomposition is $h(\mathbb P^1)=1\oplus\mathbb L$.
The tensor binomial formula gives
$40\mathbb L^{-k}(1+\mathbb L)^{2k}
 =40\sum_{a=-k}^k\binom{2k}{k+a}\mathbb L^a$;
apply AC-G04.  The construction uses finite products for each $k$.
Their cohomological weights range from $-2k$ to $2k$, which grow with
$k$.  No common finite cohomological realization or compatible
multiplication has been supplied by this calculation.
\end{proof}

Here ``actual'' is deliberately used instead of the potentially ambiguous
word ``effective'': allowing $j<0$ allows inverse Tate twists, which are
not effective motives in the customary geometric convention.  Conversely,
an actual scheme can have a virtual compactly supported cohomology class:
$\mathbb G_m$ has count $q-1$ and zeta $\zeta(s-1)/\zeta(s)$.
Thus ``no inverse zeta factors'' characterizes this specified additive
semisimple cone, not all schemes with nonnegative point counts.

\begin{proposition}[AC-G07: adjoining a point is not shifting an ideal norm]
\label{AC-G07}
Disjointly adjoining a section to a scheme over $\operatorname{Spec}\mathbb Z$
adds one to each prime-power point count and multiplies its zeta by
$\zeta(s)$.  In contrast, the norm-shifted series of the supplied note,
already for ideals of $\mathbb Z$, is
\[
 \sum_{n\geq1}(n+1)^{-s}=\zeta(s)-1\qquad(\Re s>1).
\]
It is not the Hasse--Weil zeta of $\mathbb P^1_{\mathbb Z}$,
which equals $\zeta(s)\zeta(s-1)$.
\end{proposition}
\begin{proof}
In the exponential definition of each local zeta, adding one to every
point count adds $\sum_{n\geq1}p^{-ns}/n$ to its logarithm and hence
multiplies the factor by $(1-p^{-s})^{-1}$.
The norm-shifted formula follows by replacing $n+1$ by $m\geq2$.
The projective line has count $p^n+1$, so the same exponential calculation
gives its two factors.  For example, the shifted series has zero
Dirichlet coefficient at $1$, whereas a normalized Euler product and
$\zeta(s)\zeta(s-1)$ have coefficient one.  Thus the constructions differ
before any question about zeros arises.  A's theorem on primitive
Dirichlet functions cannot be substituted for a zero theorem on
$\zeta(s)-1$ by this counting analogy.
\end{proof}

\subsection{A precise split-zero construction and its scope}

For a commutative unital ring $R$, let $G(R)=R\amalg\{\tau\}$.
Extend ring addition by making $\tau$ its identity and ring multiplication
by making $\tau$ absorbing.  This is a commutative semiring, whose zero
is $\tau$ and whose unit is $1_R$; write $e=0_R$ for the supported zero.
The elementary semiring identities are exactly those in the supplied
split-zero notes.  They must not be interchanged with statements about
monoid algebras that forget ring addition.

\begin{proposition}[AC-G08: the sign and the supported zero, corrected]
\label{AC-G08}
Multiplication by $\varepsilon=-1_R$ fixes
$\{\tau\}\cup R[2]$.  If $R$ has no nonzero $2$-torsion, this fixed set
has Boolean operations on $\{\tau,e\}$, but its inclusion in $G(R)$
is not unital when $R\ne0$.  For $R=\mathbb Z$,
$\varepsilon^{2m}=1_R$ is \emph{not} fixed by multiplication by
$\varepsilon$.  Thus the claim in the supplied split-zero note \S8.2
that even powers land in the Boolean fixed locus is false.
There is a surjective semiring map
\[
 \mathbb N[\varepsilon]/(\varepsilon^2-1)\longrightarrow G(\mathbb Z),
 \quad (a,b)\longmapsto
 \begin{cases}\tau&(a,b)=(0,0),\\a-b& a+b>0.
 \end{cases}
\]
The supported zero is the image of all positive multiples of
$1+\varepsilon$, not equality with that class in the representation
semiring.  The genuine Artin object $1\oplus\varepsilon$ has dimension
$2$ and zeta $\zeta_{\mathbb Q(i)}(s)$.
\end{proposition}
\begin{proof}
For a supported element, $-r=r$ is exactly $2r=0$; $\tau$ is fixed.
On $\{\tau,e\}$ addition is maximum and multiplication is minimum after
ordering $\tau<e$.  Its multiplicative unit is $e$, not $1_R$.
In $G(\mathbb Z)$, $\varepsilon\cdot1_R=-1_R\ne1_R$, proving the
refutation.  Addition of two nonzero coefficient pairs remains nonzero;
their product has total coefficient $(a+b)(c+d)>0$ and signed value
$(a-b)(c-d)$.  This proves the map respects both operations, including
zero.  Every integer is attained by a nonzero pair; zero signed value
occurs precisely at $(a,b)=(m,m)$, $m>0$.
In the representation semiring, $(1+\varepsilon)^2=2(1+\varepsilon)$,
whereas $e^2=e$, so the classes cannot be identified there.
The arithmetic assertion follows from the trivial/sign decomposition
of the permutation representation on $\mathbb Q(i)$ and AC-G04.
\end{proof}

Multiplication by $-1$ is an action on the underlying additive object;
it is not in general a unital semiring automorphism.  Its Boolean fixed
set is therefore not a subalgebra over the same unit.  At inert primes,
$1\oplus\varepsilon$ has zero odd Frobenius traces and trace $2$ at even
powers.  At split primes it has trace $2$ at every power.  The inert trace
zero is the relevant cancellation; it is not the zero object.

\begin{proposition}[AC-G09: sheafification on a fixed space]
\label{AC-G09}
Let $\mathcal O$ be a sheaf of unital rings on a topological space $X$.
The sectionwise presheaf $U\mapsto G(\mathcal O(U))$ is generally not a
sheaf.  Its associated sheaf has the explicit description
\[
 \mathcal G(U)=\{(V,r):V\subset U\text{ open and closed},\ r\in\mathcal O(V)\}.
\]
The bottom section has $V=\varnothing$.  Addition uses the union of
supports and the original ring sum on their intersection; multiplication
uses their intersection and the original ring product.
Its stalk at $x$ is $G(\mathcal O_x)$.  If $\mathcal O$ has no
$2$-torsion, the fixed sections for the sign over $U$ are precisely the
open-and-closed subsets of $U$, with Boolean union and intersection.
This construction changes the structure sheaf on $X$, not its underlying
topos or points.
\end{proposition}
\begin{proof}
A sheaf must take the empty set to a singleton, while
$G(\mathcal O(\varnothing))=G(0)$ has two elements.  Even after repairing
that value, disjoint nonempty opens reveal the defect: a supported
section on one component and $\tau$ on the other glue in a sheaf but
are not any section of $G(\mathcal O(U))$.
In the displayed construction, compatible local support indicators glue
to a locally constant indicator $U\to\{0,1\}$, hence an open-and-closed
$V$.  The supported ring sections glue uniquely by the sheaf property of
$\mathcal O$.  This proves descent.  The prescribed algebra operations
are defined on the disjoint open-and-closed pieces and satisfy the
semiring axioms stalkwise.  Every displayed section is locally either
$\tau$ or a section of $\mathcal O$, so it is locally in the image of
the presheaf.  Two presheaf sections with equal images have equal germs.
These facts prove the asserted associated-sheaf property and identify
the stalks.  For the sign, equality requires $2r=0$ on $V$; under the
hypothesis this means $r=0$, leaving just $V$.
\end{proof}

The relevant general topos input is SGA~4, Expos\'e II, Theorem 4.1(1):
the associated-sheaf functor is exact and commutes with colimits.
At the presheaf level, $G(A\times_C B)=G(A)\times_{G(C)}G(B)$: a compatible
pair is either $(\tau,\tau)$ or a supported pair in $A\times_C B$.
Sheafification therefore preserves this connected gluing identity.
The terminal object and disconnected covers require the correction just
proved; preservation of fibre products alone never proved them.

\begin{proposition}[AC-G10: a distinct construction that adds a generic point]
\label{AC-G10}
In the ordinary ideal spectrum of the semiring $G(R)$ the prime ideals
are exactly
\[
 \{\tau\},\qquad \{\tau\}\amalg\mathfrak p\quad
 (\mathfrak p\in\operatorname{Spec}R).
\]
Thus it is $\operatorname{Spec}R$ with a new open generic point $\xi$;
the closed complement is $V(e)=\operatorname{Spec}R$.
More generally, for a ringed space $(X,\mathcal O)$, define $X^+=X\amalg\{\xi\}$
with open sets $\varnothing$ and $U^+=U\cup\{\xi\}$ for $U\subset X$
open.  Set $\mathcal O^+(U^+)=G(\mathcal O(U))$ and give the empty open
the terminal semiring.  This is a sheaf of semirings.  For affine $X$
it is the semiring spectrum just described, with its usual localization
structure sheaf.  The stalk at $\xi$ is $\mathbb B$.
Every unital semiring map $G(R)\to A$ to a ring factors uniquely through
$p:G(R)\to R$, $p(\tau)=0$, $p(r)=r$.
Consequently its ring-valued points are the old ring-valued points;
$|G(\mathbb F_q)|=q+1$ is a cardinality of elements, not the number of
$\mathbb F_q$-points of this affine spectrum.
\end{proposition}
\begin{proof}
An ideal containing any supported $r$ contains $er=e$, and thereafter
its supported part is a ring ideal, since multiplication by $-1_R$
provides additive inverses.  An ideal with no supported element is
$\{\tau\}$.  Products of two supported elements stay supported, even
when their ring product is zero; hence $\{\tau\}$ is prime.  The other
prime conditions are exactly those for the ring ideal $\mathfrak p$.
The asserted topology follows from $D(e)=\{\xi\}$ and
$D(r)=D_R(r)\cup\{\xi\}$ for supported $r$.
For a nonempty cover by $U_i^+$, the overlaps all contain $\xi$ and the
restriction to $\{\xi\}$ records whether a section is supported or
$\tau$.  Compatibility forces that choice to agree on the whole cover.
In the supported case the ring sections glue over the $U_i$; in the
other case they are all $\tau$.  This proves the sheaf assertion.
Localization at a supported $r$ gives $G(R_r)$ by the usual fraction
equivalence relation, including $r=0$, when both sides are $\mathbb B$.
This identifies the affine structure sheaves and gives the stalk at
$\xi$; equivalently, inverting $e$ forces $e=1$ and then forces all
supported elements to equal $1$.
For a map $\phi$ to a ring, $\phi(e)+\phi(e)=\phi(e)$ gives $\phi(e)=0$.
The restriction to $R$ is therefore a ring homomorphism, and
$\phi(\tau)=0$ determines the factorization.  In particular, over
$\mathbb F_q$ there is one such $\mathbb F_q$-valued point when the
supported constants are required to be fixed, not $q+1$.  This is a
condition on the induced map $\mathbb F_q\to\mathbb F_q$ after ring
completion; it does not presume a unital ring map
$\mathbb F_q\to G(\mathbb F_q)$, which does not exist.
\end{proof}

AC-G09 and AC-G10 are two different, fully specified constructions.
The former answers a fixed-topos sheaf problem; the latter answers an
ordinary semiring-spectrum problem.  Neither identifies that spectrum
with the Connes--Consani spherical spectrum without a comparison of
categories and structure sheaves.  In the pointed-monoid version one can
forget addition and adjoin a new basepoint to a spatial monoid.  The new
generic residue is then the two-element pointed monoid $\mathbb F_1$;
calling it the Boolean \emph{semiring} requires the additive structure
used above.

\begin{proposition}[AC-G11: the power-map and Hopf type restrictions]
\label{AC-G11}
On the pointed multiplicative monoid underlying $G(R)$, the maps
$x\mapsto x^n$, with $\tau\mapsto\tau$, are monoid endomorphisms and
preserve a nontrivial base sign $\varepsilon=-1$ exactly when $n$ is odd.
They need not be semiring endomorphisms.  For $R=\mathbb Z$, already
$n=3$ fails additivity.
For an unpointed commutative monoid $M$, its spherical monoid algebra
has the group-like diagonal over $\mathbb F_1$; when $M=(R,\cdot)$ for a nonzero ring $R$, this
has no antipode because the supported zero is not invertible.
Moreover the diagonal on \emph{all} elements cannot descend to an
$\mathbb F_{1^2}$-linear coproduct after identifying the element
$[-1_R]$ with the base sign, unless that sign has collapsed.
\end{proposition}
\begin{proof}
Commutativity gives $(xy)^n=x^ny^n$.  The sign condition is
$(-1)^n=-1$.  In $G(\mathbb Z)$,
$(1+1)^3=8\ne2=1^3+1^3$, proving failure of semiring additivity.
Over $\mathbb F_1$, the coproduct sends $[m]$ to $[m]\otimes[m]$,
and the counit sends every supported monoid element to $1$.
An antipode would give an inverse to the group-like supported zero,
which is impossible in the multiplicative monoid of a nonzero ring:
$0r=0\ne1$.  After the stated signed identification, the same diagonal
would send $\varepsilon=[-1_R]$ to
$\varepsilon\otimes\varepsilon=1$, whereas base linearity requires
$\varepsilon\otimes1=\varepsilon$.  The two are unequal when the sign
is nontrivial.  Thus one must keep a separate spatial monoid, as in the
free signed construction of 2609.00299, \S2.2, or supply a different
coproduct and verify its relations.  The displayed diagonal is not such
a coproduct on the identified signed ring monoid.
\end{proof}

The statement about the diagonal is deliberately restricted to that
specified map; it does not assert that no other bialgebra structure can
exist.  The two spherical objects
$\mathbb S[(R,\cdot)_+]$ and $HG(R)$ differ: the first remembers monoid
multiplication, the second also remembers ring addition.  The source
notes' proposed extension must choose one before Frobenius and Hopf claims
can be transferred.

\begin{proposition}[AC-G12: a global realization of the trace-unimodality example]
\label{AC-G12}
Let $Y=\operatorname{Res}_{\mathbb Q(i)/\mathbb Q}\mathbb P^1$ and
$X=Y^2\times\mathbb P^2$.  The centered motive $h(X)(3)$ has graded
dimensions and inert-prime normalized trace coefficients respectively
\[
 d=(1,5,11,14,11,5,1),\qquad c=(1,1,3,2,3,1,1),
\]
in Tate degrees $j=-3,-2,\ldots,3$.  Its multiplicities are
\[
 a^+=(1,3,7,8,7,3,1),\qquad a^-=(0,2,4,6,4,2,0).
\]
At an inert good prime $p$, its first normalized trace is
$(p+p^{-1})^2(p+1+p^{-1})$.
Its complete motivic $L$-function is AC-G04 with these multiplicities,
so A proves its nontrivial zero lines.  The graded dimensions are unimodal
while the inert trace coefficients are not.
\end{proposition}
\begin{proof}
After the quadratic base extension, $Y$ becomes $\mathbb P^1\times\mathbb P^1$.
Galois acts on the two degree-two classes by exchanging them, so their
representation is $1\oplus\varepsilon$.  The degree-zero and top classes
are invariant.  The projective bundle formula and this explicit swap give
\[
 h(Y)=1\oplus(1\oplus\varepsilon)\mathbb L\oplus\mathbb L^2.
\]
This is an isomorphism of rational Chow motives, not just an equality
of traces: the split projective-bundle correspondences are equivariant
for the exchange of factors and descend by finite Galois descent;
restriction followed by transfer is multiplication by $2$, and the
rational projectors $(1\pm\sigma)/2$ separate the two Artin summands.
K\"unneth and $h(\mathbb P^2)=1\oplus\mathbb L\oplus\mathbb L^2$ now give
\[
 [h(X)(3)]=\mathbb L^{-3}
 (1+(1+\varepsilon)\mathbb L+\mathbb L^2)^2
 (1+\mathbb L+\mathbb L^2).
\]
Evaluation at $\varepsilon=1$ gives $d$; evaluation at
$\varepsilon=-1$ gives $c$; resolving the two characters gives
$a^+=(d+c)/2$ and $a^-=(d-c)/2$, exactly as displayed.
At an inert prime $\varepsilon$ has Frobenius $-1$, so the trace is
$p^{-3}(1+p^2)^2(1+p+p^2)$, which is the asserted expression.
The descent has a smooth model over $\mathbb Z[1/2]$, but at the ramified
prime we mean the inertia-invariant motivic factor of AC-G04, not an
unverified point count of a chosen bad integral model.  The elementary
inequalities in the two displayed coefficient sequences prove the last
claim: $d$ rises to its central coefficient and then falls, whereas
$c$ contains the strict dip $3,2,3$.
\end{proof}

\subsection{What is not settled by these constructions}

The exact remaining geometric gates are the trace-compatible
intersection theory and surface Riemann--Roch theorem on the relevant
square, together with an effectivity/descent theorem in that category.
A proves their numerical target and the constructions above give its
minimal analytic realization.  They give no functor from that Hilbert
space to algebraic cycles, no algebraicity of projectors, and no
polarization on a putative global arithmetic motive.  Accordingly they
do not prove the standard conjectures, a general motivic RH, or the
cubic Kummer window bound.

The supplied counting notes cite rank-three tables, the ``dimension
ring'' certificates, and the GCT runs by IDs F35, F41, F42, F45 and
P-F-E3-01.  The theorem AC-G06 applies to each coefficient family for
which the required Artin--Tate realization is actually given.  It does
not re-certify those absent GCT tables or prove their associativity.
An associative numerical based ring alone is also not a construction of
associators satisfying the pentagon in a tensor category; that extra
categorical realization is a separate gate.  No theorem about it is
inferred from A here.

For Grothendieck's frameworks, the precise established connections are:
AC-G01--G02 supply the numerical Hodge-index form modeled on
\emph{Sur une note de Mattuck--Tate}, J. reine angew. Math. 200 (1958),
208--215, \S1 (Hodge-index theorem and product-of-curves inequality
(1.8)), with the surface Riemann--Roch argument in \S2, (2.1).
AC-G04--G06 are constructions in Artin--Tate motives and their
$\ell$-adic realizations, with exterior powers and hence the usual
representation $\lambda$-ring structure on $R$; inertia invariants
explain AC-G05.  AC-G09 is a sheaf construction under SGA~4,
Expos\'e II, Theorem 4.1; AC-G10 is a semiring analogue and is not an
EGA scheme since its affine coefficient objects are not rings.
The RR duality paper explicitly compares its Theorem 5.3 with the
Grothendieck--Serre letter dated 15 December 1955; that comparison is
not a new theorem identifying the categories.

The Montpellier catalogue verifies the relevant archival subjects:
cote 2 (Riemann--Roch notes), cote 10 (tensor categories), cote 11
(correspondences), cote 14 (Hodge--Tate--Lefschetz--Weil cycles), cote 16
(correspondences and Poincar\'e--Hodge algebras), and cote 19 (topos).
No theorem-level scan from those cotes was successfully retrieved in
this workstream.  These are catalogue mappings only.  An exact
cote-and-scan attribution remains an explicit bibliography gate; it
must not be manufactured from the subjects of the folders.
